% explainer-source-hash: 84b3faa7d375074fddbf34456fc0a24bd6fd8f43a93e709f80ad047a903d6daa
% explainer-source-commit: cf47a30132e1d7fe8e21b690e97c240372f801d5
% explainer-generated: 2026-07-16
% Regenerate with /explain-all (see WIRING.md). Do not edit this stamp by hand:
% it is how the toolchain knows whether this document still matches the code.
\documentclass[11pt,a4paper]{article}
\input{preamble}

\title{\vspace{-1.4cm}Blog Posting Pipeline: Demo\\[4pt]
\large Brief}
\author{Explainer for \file{blog-posting-pipeline-demo}}
\date{}

\begin{document}
\maketitle
\thispagestyle{fancy}

\section*{The short version}

\begin{honest}{What this is}
A static web page that \emph{acts out} a private production system rather than being one.
The real \file{blog-posting-pipeline} writes and publishes marketing articles for paying
clients; its source stays private because the pipeline is the agency's product. This
repository is its public stand-in: three hand-written example runs for an invented coffee
brand, replayed on a timer, so a reader can see the real architecture's shape without the
real system being exposed.

Nothing it ``generates'' was generated. No model is called, no request is made. Every run
number was typed into a JavaScript file by hand.

One thing on the page \emph{is} real, and that is the point: the Blog Library holds fifteen
genuine articles the real pipeline actually published for a real client. They are the only
evidence on the page. The runs beside them are not evidence, and are labelled as such.
\end{honest}

\section{Why it exists}

``Trust me, I built a seven-stage AI pipeline'' is not evidence, and the source cannot be
shown. So the project asks a narrower question: can the \emph{architecture} be demonstrated
honestly, in a form someone can click, while the implementation stays private?

The trade it makes: give up all of the real system's substance (prompts, client rules,
code) and keep only its shape (seven stages, in order, with the right kind of data moving
between them). Then be relentless about labelling, because a convincing demo is exactly the
thing most likely to be mistaken for the real product.

\section{The provenance ladder}

The central design idea. Three kinds of content share one page, and confusing them would be
the failure. Each tier gets its own disclosure mechanism.

\begin{figure}[H]
\centering
\resizebox{\textwidth}{!}{
\begin{tikzpicture}[node distance=6mm]
  \node[box, fill=warnred!12, draw=warnred, text width=45mm, minimum height=15mm] (t1)
    {\textbf{Fabricated}\\[2pt]\footnotesize 3 example runs for an\\ invented brand};
  \node[box, text width=45mm, minimum height=15mm, right=9mm of t1] (t2)
    {\textbf{Real, third party}\\[2pt]\footnotesize 3 open-licensed stock\\ photos};
  \node[box, fill=okgreen!12, draw=okgreen, text width=45mm, minimum height=15mm, right=9mm of t2] (t3)
    {\textbf{Real, the client's}\\[2pt]\footnotesize 15 published articles the\\ pipeline produced};

  \node[extbox, below=8mm of t1, text width=45mm] (m1)
    {\footnotesize pinned banner, card ribbon,\\ ``(fictional)'' in the brand name};
  \node[extbox, below=8mm of t2, text width=45mm] (m2)
    {\footnotesize caption: ``real stock photo,\\ not AI-generated'' + CREDITS};
  \node[extbox, below=8mm of t3, text width=45mm] (m3)
    {\footnotesize provenance footer linking to\\ the genuine original URL};

  \draw[flow] (t1) -- (m1); \draw[flow] (t2) -- (m2); \draw[flow] (t3) -- (m3);
\end{tikzpicture}}
\caption{Three origins, three disclosure mechanisms.}
\end{figure}

The tempting design is to invent a fake blog library too, so everything is uniformly
fictional and no disclosure question arises. The project rejects that: the client's posts
are already public, the real pipeline genuinely produced them, so showing them (with a link
back to each original) is the only actual proof on the page that the system ships anything.
The cost of that choice is that real and fake now sit in one grid, and the labelling carries
the whole burden.

\section{Architecture at a glance}

No server, no framework, no build step, no dependencies. About 1,700 lines of hand-written
code across four files.

\begin{figure}[H]
\centering
\resizebox{0.93\textwidth}{!}{
\begin{tikzpicture}[node distance=8mm]
  \node[box, text width=30mm] (html) {\file{index.html}\\[2pt]\scriptsize empty panels,\\ banner, tabs};
  \node[store, right=15mm of html, text width=30mm] (data) {\file{demo-data.js}\\[2pt]\scriptsize content only,\\ 3 constants};
  \node[box, right=15mm of data, text width=30mm] (app) {\file{app.js}\\[2pt]\scriptsize logic only,\\ renderers + state};
  \node[extbox, above=8mm of data, text width=32mm] (css) {\file{styles.css}\\[2pt]\scriptsize design tokens from\\ the real client site};

  \draw[flow] (html) -- node[lbl, above] {loads 1st} (data);
  \draw[flow] (data) -- node[lbl, above] {loads 2nd} (app);
  \draw[dflow] (css) -- (html);

  \node[box, below=13mm of html, text width=31mm] (p1) {Panel 1: run};
  \node[box, below=13mm of data, text width=31mm] (p2) {Panel 2: library};
  \node[box, below=13mm of app, text width=31mm] (p3) {Panel 3: reader};
  \draw[flow] (app.south) ++(0,0) -- ++(0,-4mm) -| (p3.north);
  \draw[flow] (app.south) ++(0,0) -- ++(0,-4mm) -| (p2.north);
  \draw[flow] (app.south) ++(0,0) -- ++(0,-4mm) -| (p1.north);
  \node[lbl, right=2mm of app, text width=26mm] {no \code{fetch},\\ no \code{XHR},\\ anywhere};
\end{tikzpicture}}
\caption{The clean split: \file{demo-data.js} is content with no behaviour,
\file{app.js} is behaviour with no content.}
\end{figure}

The data file holds three constants: \code{REAL\_DAVONEX\_POSTS} (15 real posts, full
article bodies, roughly 154,000 characters, which is most of the file's 231 KB),
\code{DEMO\_RUNS} (3 fabricated runs of 7 steps each), and \code{PRESETS} (the chip labels,
which also link each fabricated run to the real post whose shape inspired it).

\section{The seven stages}

The rail across the top is the demo's central claim: \emph{this is the real system's
architecture}.

\begin{figure}[H]
\centering
\resizebox{\textwidth}{!}{
\begin{tikzpicture}[node distance=4mm]
  \node[box, text width=20mm, minimum height=11mm] (s1) {\textbf{1. Brief}};
  \node[box, text width=20mm, minimum height=11mm, right=5mm of s1] (s2) {\textbf{2. Keywords}};
  \node[box, text width=20mm, minimum height=11mm, right=5mm of s2] (s3) {\textbf{3. Section plan}};
  \node[box, text width=20mm, minimum height=11mm, right=5mm of s3] (s4) {\textbf{4. Writer}\\[1pt]\scriptsize agentic};
  \node[box, text width=20mm, minimum height=11mm, right=5mm of s4] (s5) {\textbf{5. Image prompt}};
  \node[box, text width=20mm, minimum height=11mm, right=5mm of s5] (s6) {\textbf{6. Image}};
  \node[box, text width=20mm, minimum height=11mm, right=5mm of s6] (s7) {\textbf{7. Publish}};
  \draw[flow] (s1) -- (s2); \draw[flow] (s2) -- (s3); \draw[flow] (s3) -- (s4);
  \draw[flow] (s4) -- (s5); \draw[flow] (s5) -- (s6); \draw[flow] (s6) -- (s7);
  \node[extbox, above=8mm of s2, text width=28mm] (res) {\scriptsize AEO/SEO research,\\ unnumbered stage};
  \draw[dflow] (s1.north) |- (res.west);
  \draw[dflow] (res.east) -| (s2.north);
  \node[lbl, below=2mm of s1] {pro}; \node[lbl, below=2mm of s2] {pro};
  \node[lbl, below=2mm of s3] {flash}; \node[lbl, below=2mm of s4] {pro};
  \node[lbl, below=2mm of s5] {flash}; \node[lbl, below=2mm of s6] {image};
  \node[lbl, below=2mm of s7] {no model};
\end{tikzpicture}}
\caption{The seven stages. The unnumbered research stage is not on the visible rail but does
appear as a line in the cost table, which is a creditable piece of fidelity.}
\end{figure}

\begin{keyidea}{Architectural fidelity: verified, and it holds}
The demo's stages were compared against the real private project's own documentation. The
stage names, their order, the tiering of expensive against cheap models, the specific model
identifiers, the fact that stage 4 is an agentic loop with exactly two tools, the unnumbered
research stage, and the publish stage's trio of actions all match.

This is the demo's strongest result: its architectural claim is accurate, not marketing.
\end{keyidea}

\section{How a ``run'' actually works}

Clicking ``Run the pipeline'' computes nothing. All seven steps are already in the page
before the click; the button reveals one more of them every 900 milliseconds by adding a CSS
class, then adds a fabricated card to the library with a fifteen-minute countdown.

Two consequences worth stating out loud:

\begin{itemize}
  \item \textbf{The run cannot fail}, because nothing runs. The real system's failure modes
    (a refusal, a timeout, a malformed reply, a rejected publish) have no representation
    here at all.
  \item \textbf{The finished article is in the page from the start.} Anyone who opens
    developer tools before clicking sees the whole post. That is the literal meaning of
    ``scripted.''
\end{itemize}

\section{What was verified, by running it}

The page was loaded in a real headless browser, driven through a full run, and
instrumented. The claims below were tested, not trusted.

\begin{center}
\begin{tabular}{@{}p{68mm}p{24mm}p{42mm}@{}}
\toprule
\textbf{Claim} & \textbf{Result} & \textbf{Evidence} \\
\midrule
``No live API calls happen here'' & \textbf{True} & Zero network requests to any non-local
origin, across the entire flow. \\
Self-contained (fonts, images local) & \textbf{True} & The only \code{url()} in the CSS is
a bundled font. No CDN. \\
Free of JavaScript errors & \textbf{True} & Zero page and console errors. \\
Seven stages step through & \textbf{True} & 0 of 7 visible before the click, 7 of 7 after. \\
``All 15 real posts'' & \textbf{True} & 15 entries, 15 image files, no broken paths, no
orphans. \\
The countdown and expiry work & \textbf{True} & Counter went (15) to (16); ribbon ticked
14:58 to 14:56. \\
Cost range is the real published figure & \textbf{Partly false} & Only 1 of 3 presets
matches. See below. \\
\bottomrule
\end{tabular}
\end{center}

\begin{keyidea}{The headline result}
The demo's central promise, that it is inert and self-contained, is true and provably so.
Driven through its whole flow in a real browser, the page issued \textbf{no network requests
whatsoever}.
\end{keyidea}

\section{The findings that matter}

\begin{honest}{The cost caption claims fabricated numbers are real}
Under every run, the page prints: ``The total range is the real, published per-post cost for
the production system (see its README).'' That is a claim of fact, and it is true for only
one of the three presets. The demo gives each preset a different total range; checked against
the real project's own published figure, the first matches exactly and the other two do not.
They overshoot it at both ends. They are fabricated variations displayed under a caption
asserting they are real.

The demo's README compounds this by describing its headline figure as ``anchored to'' the
source project's published range. The number it quotes is the spread across the demo's own
three presets, not the source's published range. It reads as a citation and is not one.

This is the one place where a project built entirely around honest labelling labels
fabricated content as real. It is fixable in minutes: set all three presets to the real
published range, or say the totals are illustrative too. The first is better. Fix it before
showing the page to anyone.
\end{honest}

\begin{honest}{The library grid is rebuilt from scratch every second, for 15 minutes}
The countdown ribbon ticks via \code{setInterval(renderLibrary, 1000)}, but
\code{renderLibrary} does not update the ribbon: it reassigns the whole grid's
\code{innerHTML}, destroying and recreating all 16 cards and all 16 images, roughly 900
times over a card's lifetime, to change five characters.

Measured, not inferred: a mutation observer recorded 4 grid rebuilds in 4 seconds with a
demo card live, and 0 without one. It is observable, too. An attempt to click a card through
browser automation failed with ``Node is detached from document,'' because the card was
destroyed between being found and being clicked. Real users survive it only because the click
handler is bound to the container rather than the cards; what does not survive is text
selection, which is wiped once a second. The timer also keeps rebuilding the grid while it is
hidden behind the reader panel.
\end{honest}

\begin{honest}{The articles claim citations they do not contain}
The README describes the Writer step as producing ``numbered citations with a References
list,'' matched deliberately to the real system's prompt templates. The reference lists are
real and numbered; the inline citations are absent. Searching all three fabricated articles
for a marker of the form \code{[1]} returns zero matches. The prose cites nothing and a
numbered source list simply appears at the end. Since the entire editorial argument for this
kind of content is that claims are attributable, this is a gap in the demo's most
load-bearing artifact.
\end{honest}

Smaller, in descending order of interest:

\begin{itemize}
  \item \textbf{One preset's arithmetic does not close.} Its per-step line items sum to
    \$0.422, above the \$0.42 upper bound of the total printed directly below them.
  \item \textbf{No tab looks selected while reading a post.} The reader is a third panel with
    no tab button, so \code{switchTab("reader")} matches nothing and deactivates both tabs.
    Confirmed in the browser: zero tabs marked active while reading.
  \item \textbf{\code{escapeHtml} does not escape quotes} yet its output is dropped into HTML
    attributes. Not exploitable today, since every string is a hand-written constant and no
    user input reaches a renderer, but it is a guard that looks complete and is not.
  \item \textbf{The same sources render two ways:} step 4 prints explicit \code{[1]},
    \code{[2]} markers from a \code{ref} field, while the step 7 preview drops \code{ref} and
    lets the list auto-number. They agree today only by luck of ordering.
  \item \textbf{All six fabricated sources point at \code{example.com}.} The honest choice,
    but the References lists lead nowhere.
  \item \textbf{Post cards are not keyboard accessible:} plain \code{<article>} elements with
    no \code{tabindex}, opened by a delegated click.
  \item \textbf{The page is \code{noindex}}, which is correct: it reproduces a client's
    articles verbatim and must not compete with the client's own pages for the rankings the
    real pipeline exists to win.
\end{itemize}

\section{Licensing: handled unusually well}

Three licences apply to one folder. The repository's own code is PolyForm Strict. Two stock
photos are public domain and CC BY. The third is CC BY-SA 4.0, and \file{CREDITS.md}
explicitly records that this one cropped file stays CC BY-SA and does \emph{not} fall under
the repository's licence.

\begin{keyidea}{The share-alike catch, spotted}
CC BY-SA is viral: a derivative must carry the same licence, and a cropped photo is a
derivative. Dropping it into a strictly licensed repository without comment would quietly
violate its terms. Most portfolio projects never notice. This one carved it out by name, and
can say so.
\end{keyidea}

The fifteen real posts are not open-licensed at all; they are the client's. They are
reproduced with a provenance link on every article and a \file{CREDITS.md} explaining the
reasoning. The demo also refuses to hotlink: everything is fetched once and stored locally,
which is what makes the ``no live calls'' guarantee airtight rather than nearly airtight.

\section{What it cannot demonstrate}

\begin{figure}[H]
\centering
\resizebox{0.97\textwidth}{!}{
\begin{tikzpicture}[node distance=6mm]
  \node[box, fill=okgreen!10, draw=okgreen, text width=58mm, minimum height=34mm, align=left] (can)
    {\textbf{Does show}\\[3pt]
     \footnotesize $\bullet$ The real stage list, order, model\\ \quad tiering (independently verified)\\[2pt]
     $\bullet$ The shape of the data and the output\\[2pt]
     $\bullet$ That the real system ships: 15 real posts\\[2pt]
     $\bullet$ Front-end craft and disclosure design\\[2pt]
     $\bullet$ Judgment about what may be shown};
  \node[box, fill=warnred!10, draw=warnred, text width=58mm, minimum height=34mm, align=left, right=9mm of can] (cannot)
    {\textbf{Cannot show}\\[3pt]
     \footnotesize $\bullet$ Any prompt or prompt engineering\\[2pt]
     $\bullet$ That the pipeline works (nothing runs)\\[2pt]
     $\bullet$ Output quality (hand-written to look good)\\[2pt]
     $\bullet$ Failure handling (no error path exists)\\[2pt]
     $\bullet$ Latency (900 ms is a UI pause)\\[2pt]
     $\bullet$ Cost measurement (typed in; 2 are wrong)\\[2pt]
     $\bullet$ Any backend skill at all};
  \draw[flow, <->] (can) -- (cannot);
\end{tikzpicture}}
\caption{The boundary of the evidence.}
\end{figure}

\begin{honest}{How to talk about it}
The worst possible framing is ``here is my AI pipeline running.'' It is not running, one
question exposes that, and every other claim then becomes suspect.

The strong framing is: ``the real system is private, so I built an honest replica of its
architecture and put it beside fifteen posts the real system actually published. The runs are
fabricated and labelled as such. Here is what it proves, and here is what it cannot.'' That
turns the demo's greatest weakness, that it is fake, into evidence of the judgment that
produced it. Fix the cost caption first: it is the one place the page's own discipline slips,
and it slips in the exact direction an interviewer will probe.
\end{honest}

\section*{Glossary}

\begin{description}[leftmargin=0pt, style=nextline, itemsep=3pt]
  \item[Static site] A site made only of files handed over unchanged, with no program running
    on a server. Everything happens in your browser. This is why the demo can make no live
    calls: there is nothing to call.
  \item[Vanilla JavaScript] JavaScript with no framework. The programmer edits the page
    directly. Removes a whole toolchain; costs you the free optimizations a framework gives,
    one of which is missed here.
  \item[DOM] The browser's live, tree-shaped model of the page. JavaScript changes what you
    see by editing this tree.
  \item[innerHTML] Assign to it and the browser throws away an element's contents and rebuilds
    them from scratch. Quick to write; it destroys rather than updates, which is the root of
    the once-a-second rebuild defect.
  \item[JSON] A plain-text format for structured data. How the real system's stages hand
    results to each other; showing it raw signals ``a machine read this, not a human.''
  \item[SEO / AEO] SEO shapes a page so a search engine ranks it and a human clicks. AEO
    shapes a page so an AI assistant answering the question directly cites you as its source.
    With AEO there may be no click at all, so the goal shifts from ranking to being quotable.
  \item[Agentic loop] The model is given tools and loops (think, call a tool, read the result,
    think again) until it decides it is done, rather than answering once. Stage 4 does this
    with two tools, which is why it dominates the cost table.
  \item[Embedding / Pinecone] An embedding turns text into a list of numbers placed so similar
    meanings sit near each other; Pinecone searches those by nearness, not keyword. The real
    system uses it to avoid repeating itself and to find internal links.
  \item[IndexNow] A protocol for telling search engines ``this URL just changed, come look
    now.'' One request. Shown here as \code{not sent (demo)}.
  \item[Slug] The readable identifier in a web address. Here it doubles as each post's key.
  \item[Design tokens] Named values for a design's raw ingredients (a colour, a corner
    radius), defined once and referenced everywhere so the design stays consistent.
  \item[Escaping] Rewriting characters like \code{<} into harmless stand-ins so text is
    displayed rather than executed. Skipping it causes cross-site scripting bugs.
  \item[PolyForm Strict] A licence granting the right to read and evaluate, reserving
    essentially everything else. Standard for showing work without licensing its reuse.
  \item[CC BY-SA] An open licence with a share-alike term: derivatives must carry the same
    licence. Viral, and easy to violate by accident.
\end{description}

\vfill
\begin{center}\footnotesize\itshape
Every finding here was taken from the source, or from driving the page in a real browser.
Claims about the private \file{blog-posting-pipeline} were checked against that project's own
documentation; none of its private detail is reproduced here.
\end{center}

\end{document}
